\chapter{Methodology}
\label{chap:methodology}

\section{Methodological Architecture Overview}
The IntelliTwin methodological framework links machine learning degradation modeling, distribution-free uncertainty calibration, and sequential decision simulation into a cohesive prognostic pipeline. The pipeline operates under three immutable architectural principles:
\begin{enumerate}
    \item \textbf{Causal Temporal Isolation}: Predictions condition strictly on past operational history within a 30-cycle causal horizon; zero-padding and future lookahead are strictly precluded.
    \item \textbf{Engine-Level Disjoint Calibration}: Conformal nonconformity scores are computed exclusively on dedicated, reserved calibration engine trajectories that share no engine identities with training or held-out test splits.
    \item \textbf{Asymmetric Cost Alignment}: Nonconformity metrics penalize dangerous RUL overestimation more severely than premature maintenance indications.
\end{enumerate}

\section{Data Preprocessing and Causal Windowing}
\subsection{Feature Space Construction}
From the raw C-MAPSS telemetry streams, seven invariant sensor channels showing zero or near-zero variance across operating regimes are permanently removed during training-only qualification:
\begin{itemize}
    \item Excluded sensors: $s_1$ (total temperature at fan inlet), $s_5$ (total pressure at LPC outlet), $s_6$ (total pressure at HPC outlet), $s_{10}$ (engine pressure ratio), $s_{16}$ (bypass ratio), $s_{18}$ (demanded fan speed), and $s_{19}$ (demanded corrected fan speed).
\end{itemize}
The retained 14 sensor signals alongside 3 operational settings constitute the $17$-dimensional cycle state vector $\mathbf{x}_{i,t} \in \mathbb{R}^{17}$. Min-max scaling parameters $\boldsymbol{\mu}_{\text{min}}, \boldsymbol{\mu}_{\text{max}} \in \mathbb{R}^{17}$ are fitted strictly on training engine splits and applied out-of-sample to calibration, validation, and test splits:
\begin{equation}
\tilde{\mathbf{x}}_{i,t} = \frac{\mathbf{x}_{i,t} - \boldsymbol{\mu}_{\text{min}}}{\boldsymbol{\mu}_{\text{max}} - \boldsymbol{\mu}_{\text{min}}}.
\end{equation}

\subsection{Causal Windowing Protocol}
Temporal feature matrices are formed via rolling causal windows of duration $w = 30$ operational cycles with stride $s = 1$:
\begin{equation}
\mathbf{W}_{i,t} = [\tilde{\mathbf{x}}_{i, t-29}, \tilde{\mathbf{x}}_{i, t-28}, \dots, \tilde{\mathbf{x}}_{i,t}] \in \mathbb{R}^{30 \times 17}, \quad t \ge 30.
\end{equation}
To maintain rigorous physical validity, windows are formed exclusively within continuous engine trajectories; windows never span across engine boundaries, and cycles $t < 30$ are unpredicted.

\section{Base Degradation Predictors}
We examine two distinct algorithmic paradigms receiving identical 30-cycle information horizons:
\begin{enumerate}
    \item \textbf{Gradient Boosted Regressor (GBR)}: A LightGBM tree ensemble operating on the flattened feature vector:
    \begin{equation}
    \mathbf{v}_{i,t} = \operatorname{vec}(\mathbf{W}_{i,t}) \in \mathbb{R}^{510}.
    \end{equation}
    Hyperparameters (tree depth, learning rate, minimum leaf split) were selected during Phase 5 validation.
    \item \textbf{Long Short-Term Memory (LSTM)}: A two-layer recurrent neural network architecture processing the temporal sequence $\mathbf{W}_{i,t} \in \mathbb{R}^{30 \times 17}$. Hidden state dynamics capture progressive degradation accumulation over time.
\end{enumerate}

To ensure complete empirical replicability and isolate model training variance, five random seeds ($2026, 2027, 2028, 2029, 2030$) are executed for each architecture across all three C-MAPSS sub-datasets (FD001, FD002, FD004), yielding 30 frozen base model checkpoints in the registry.

\section{Conformal Calibration Formulation}
\subsection{Engine-Weighted Quantile Computation}
Standard split-conformal inference pools calibration residuals across all observation cycles, implicitly assigning greater statistical weight to longer-lived engines. To prevent survivor bias and enforce equal contribution across asset units, IntelliTwin defines an exact engine-weighted empirical distribution.

Let $\mathcal{D}_{\text{cal}} = \{(\mathbf{W}_{i,t}, y_{i,t}) : i \in \mathcal{E}_{\text{cal}}, t \in [w, T_i]\}$ denote the reserved calibration dataset containing $|\mathcal{E}_{\text{cal}}| = M$ engines. Each observation $(i,t)$ receives a normalized rational weight:
\begin{equation}
w_{i,t} = \frac{1}{M \cdot (T_i - w + 1)}.
\end{equation}
The empirical cumulative distribution function of nonconformity scores $\alpha_{i,t}$ is given by:
\begin{equation}
\hat{F}_{\text{cal}}(s) = \sum_{(i,t) \in \mathcal{D}_{\text{cal}}} w_{i,t} \cdot \mathbf{1}_{\{\alpha_{i,t} \le s\}}.
\end{equation}
For a specified significance level $\alpha \in (0, 1)$, the conformal quantile $q^*$ is computed as:
\begin{equation}
q^*(\alpha) = \inf \left\{ s \in \mathbb{R} : \hat{F}_{\text{cal}}(s) \ge \frac{\lceil (M+1)(1-\alpha) \rceil}{M} \right\}.
\end{equation}

\subsection{Conformal Calibration Variants}
We formalize three conformal calibration variants:
\begin{enumerate}
    \item \textbf{Global HCP (\texttt{global\_hcp})}: Evaluates absolute residual nonconformity:
    \begin{equation}
    \alpha_{i,t} = |y_{i,t} - \hat{y}_{i,t}|.
    \end{equation}
    The resulting prediction interval on test engine $j$ is symmetric:
    \begin{equation}
    \mathcal{C}_{\text{global}}(\mathbf{W}_{j,t}) = [\hat{y}_{j,t} - q^*, \hat{y}_{j,t} + q^*].
    \end{equation}

    \item \textbf{Regime-Stratified HCP (\texttt{regime\_hcp})}: Conditions nonconformity scoring on operating regime cluster $k(t) \in \{1, \dots, 6\}$:
    \begin{equation}
    \alpha_{i,t}^{(k)} = |y_{i,t} - \hat{y}_{i,t}| \quad \text{for } \mathbf{x}_{i,t} \in \mathcal{R}_k.
    \end{equation}
    To prevent high quantile variance in sparse flight conditions, if regime $k$ contains fewer than $N_{\text{min}} = 100$ calibration observations, the framework executes a deterministic fallback to the pooled global nonconformity quantile $q^*_{\text{global}}$.

    \item \textbf{Decision-Aware HCP (\texttt{decision\_aware\_hcp})}: Incorporates asymmetric penalty scaling governed by parameter $\beta = -\ln 2$:
    \begin{equation}
    \alpha_{i,t} = \begin{cases}
    \frac{\hat{y}_{i,t} - y_{i,t}}{a(p_{i,t})} & \text{if } y_{i,t} \le \hat{y}_{i,t} \text{ (overestimation)}, \\
    \frac{y_{i,t} - \hat{y}_{i,t}}{b(p_{i,t})} & \text{if } y_{i,t} > \hat{y}_{i,t} \text{ (underestimation)},
    \end{cases}
    \end{equation}
    where $p_{i,t} = \sigma((\tau_{\text{crit}} - \hat{y}_{i,t})/\gamma)$ represents the estimated cycle failure probability, and scaling factors satisfy:
    \begin{equation}
    a(p_{i,t}) = \exp(-\beta \cdot p_{i,t}), \quad b(p_{i,t}) = \exp(\beta \cdot p_{i,t}).
    \end{equation}
    The resulting asymmetric prediction interval on test observations is:
    \begin{equation}
    \mathcal{C}_{\text{asym}}(\mathbf{W}_{j,t}) = [\hat{y}_{j,t} - q^* \cdot a(p_{j,t}), \hat{y}_{j,t} + q^* \cdot b(p_{j,t})].
    \end{equation}
\end{enumerate}

\section{Discrete-Event Maintenance Simulation}
The sequential maintenance policy executes step-by-step across unseen test engine trajectories. At cycle $t$, if the lower prediction bound satisfies $L_{j,t}(\alpha) \le \tau_{\text{lead}}$, a scheduled maintenance order is issued. The engine operates for $\tau_{\text{lead}} = 2$ cycles to simulate logistical dispatch, and is overhauled at $t + \tau_{\text{lead}}$. If the engine reaches its true physical failure cycle $T_j$ before the scheduled overhaul completes, an in-flight catastrophic failure is recorded.
